% ============================================================
% RIHB-in-the-loop, v3.
% Two contributions: live on-device SAR receipts (Part 1) +
% differentiable-pose RIHB capacity (Part 2), unified by the
% same body-twin machinery v5 already deploys at the BS side.
% Author: R. Wydaeghe
% ============================================================

\documentclass[11pt,a4paper]{article}

\usepackage[margin=2.4cm]{geometry}
\usepackage[utf8]{inputenc}
\usepackage[T1]{fontenc}
\usepackage{lmodern}
\usepackage{microtype}
\linespread{1.04}

\usepackage{amsmath,amssymb,amsthm,mathtools}
\usepackage{bm}
\usepackage{mathrsfs}
\usepackage{bbm}
\usepackage{cancel}
\usepackage{xfrac}
\usepackage{booktabs}
\allowdisplaybreaks[2]

\usepackage{empheq}
\usepackage[most]{tcolorbox}
\tcbset{highlight math style={%
    enhanced, colframe=black!70, colback=gray!6,
    arc=2pt, boxrule=0.5pt, left=3pt, right=3pt}}

\usepackage{siunitx}
\sisetup{detect-all, range-phrase = {--}, range-units = single,
         per-mode = symbol, inter-unit-product = {\,}}

\usepackage{enumitem}
\usepackage{xcolor}
\usepackage{xspace}
\usepackage{fancyhdr}
\pagestyle{fancy}
\fancyhf{}
\cfoot{\thepage}
\renewcommand{\headrulewidth}{0pt}

\usepackage{titlesec}
\titleformat{\section}{\large\bfseries}{\thesection}{1em}{}
\titleformat{\subsection}{\normalsize\bfseries}{\thesubsection}{1em}{}

\usepackage{tikz}
\usetikzlibrary{positioning,arrows.meta,calc,fit,backgrounds,shapes.geometric}

\usepackage[colorlinks=true, linkcolor={blue!55!black},
            citecolor={green!45!black}, urlcolor={blue!65!black}]{hyperref}
\usepackage[nameinlink,capitalize]{cleveref}

\usepackage{aliascnt}
\theoremstyle{definition}
\newtheorem{theorem}{Theorem}[section]
\newaliascnt{proposition}{theorem}
\newtheorem{proposition}[proposition]{Proposition}
\aliascntresetthe{proposition}
\newaliascnt{lemma}{theorem}
\newtheorem{lemma}[lemma]{Lemma}
\aliascntresetthe{lemma}
\newaliascnt{corollary}{theorem}
\newtheorem{corollary}[corollary]{Corollary}
\aliascntresetthe{corollary}
\newaliascnt{definition}{theorem}
\newtheorem{definition}[definition]{Definition}
\aliascntresetthe{definition}
\newaliascnt{remark}{theorem}
\newtheorem{remark}[remark]{Remark}
\aliascntresetthe{remark}
\crefname{proposition}{Proposition}{Propositions}
\Crefname{proposition}{Proposition}{Propositions}
\crefname{lemma}{Lemma}{Lemmas}\Crefname{lemma}{Lemma}{Lemmas}
\crefname{corollary}{Corollary}{Corollaries}\Crefname{corollary}{Corollary}{Corollaries}
\crefname{definition}{Definition}{Definitions}\Crefname{definition}{Definition}{Definitions}
\crefname{remark}{Remark}{Remarks}\Crefname{remark}{Remark}{Remarks}

\newcommand{\khat}{\hat{\bm{k}}}
\newcommand{\nhat}{\hat{\bm{n}}}
\newcommand{\rhat}{\hat{\bm{r}}}
\newcommand{\rr}{\mathbf{r}}
\newcommand{\EE}{\mathbf{E}}
\newcommand{\Sinc}{S_{\mathrm{inc}}}
\newcommand{\Sab}{S_{\mathrm{ab}}}
\DeclareMathOperator{\ReLU}{ReLU}
\DeclareMathOperator{\RE}{Re}
\DeclareMathOperator{\IM}{Im}
\DeclareMathOperator{\tr}{tr}
\DeclareMathOperator*{\argmin}{arg\,min}
\DeclareMathOperator*{\argmax}{arg\,max}
\newcommand{\diff}{\mathrm{d}}
\newcommand{\ident}{\mathbbm{1}}
\newcommand{\Reals}{\mathbb{R}}
\newcommand{\Complex}{\mathbb{C}}
\newcommand{\Gtil}{\tilde{\mathbf{G}}}
\newcommand{\Qabs}{\mathbf{Q}_{\mathrm{ab}}}
\newcommand{\Frib}{\mathbf{F}_{\mathrm{b}}}
\newcommand{\Phimat}{\mathbf{\Phi}}
\newcommand{\Mmat}{\mathbf{M}}
\newcommand{\Jmat}{\mathbf{J}}
\newcommand{\hh}{\mathbf{h}}
\newcommand{\xvec}{\mathbf{x}}
\newcommand{\thetab}{\bm{\theta}}
\newcommand{\betab}{\bm{\beta}}
\newcommand{\Pabs}{P_{\mathrm{ab}}}
\newcommand{\SINR}{\mathrm{SINR}}
\newcommand{\kout}{\hat{\bm{k}}_{\mathrm{o}}}
\newcommand{\kin}{\hat{\bm{k}}_{\mathrm{i}}}

\title{\textbf{The User as a Reflective Intelligent Body:\\
Live SAR Receipts and Pose-Differentiable\\
mmWave Capacity (v3)}}
\author{Robin Wydaeghe\\
 Ghent University \& imec\\
 \texttt{robin.wydaeghe@ugent.be}}
\date{\today}

\makeatletter
\renewcommand{\maketitle}{%
  \thispagestyle{plain}%
  \begin{center}
    \rule{\linewidth}{1pt}\\[0.4em]
    {\LARGE\bfseries \@title}\\[0.5em]
    {\large \@author}\\[0.3em]
    {\small \@date}\\[0.4em]\rule{\linewidth}{1pt}
  \end{center}
  \vspace{0.5em}
}
\makeatother

\begin{document}
\maketitle

\begin{abstract}
A user with a smartphone at mmWave frequencies is, geometrically, a
\emph{reflective intelligent body} (RIB): a passive scatterer in the
propagation graph whose surface phase pattern is fixed by pose, and
whose only tuning knob is gait.  This note develops two contributions
that share a single piece of machinery (the on-device SMPL-X body twin
of \texttt{rib.tex} and \texttt{paper\_jsac\_v5.tex}) and that
together address the two reasons mmWave deployment slowed: blockage
and consumer worry.  \textbf{Part~1, the live SAR receipt}, runs the
body twin on the user's phone and computes the user's own absorbed
power in real time from the BS-shared path dictionary.  The receipt
is displayed for transparency, optionally logged for personal use,
and (with consent) aggregated across users for population-scale
exposure epidemiology.  This is the trust artifact mmWave has been
missing.  \textbf{Part~2, the RIHB capacity lift} (Reflective
Intelligent Human Body: a RIB whose tuning controller is the human
brain), uses the same body twin to compute a pose gradient on
achievable rate.  The BS suggests sub-degree to few-degree pose moves
that improve the user's channel via line-of-sight un-shadowing and
via favorable bistatic backscatter.  Both contributions are
complementary to the BS-side compliance pipeline of v5
(\texttt{paper\_jsac\_v5.tex}); v5 enforces, this paper informs and
optimizes.  The headline empirical claim is in the companion
experiment: at a binding-regime body-shadowed line-of-sight, a
\SI{15}{\degree} torso turn buys \SI{14}{\percent} of rate at no
operationally-relevant exposure cost.  The dose receipt confirms
this; nothing in the user's optimization trades it.
\end{abstract}

\vspace{0.3em}
\noindent\textit{Style note.}~Wandering companion to
\texttt{rib.tex} and \texttt{q\_complement.tex}; theorems where the
proof is short, propositions where I want a reviewer to take or
leave a sketch, conjectures flagged.  Intended home for the
deployable parts is \S\,IV-V of a JSAC2 SI submission.

\tableofcontents
\bigskip

% ============================================================
\section{Introduction: trust and capacity, both gated by the body}\label{sec:intro}
% ============================================================

mmWave was supposed to be the headline 5G technology.  Two reasons
the deployment slowed.  First, blockage: a body in the line of sight
of an FR2 channel costs \SI{25}{} to \SI{40}{dB}, and trees, walls,
and human torsos are everywhere bodies are.  Second, consumer worry:
the post-deployment communications-public-relations conversation
has been dominated by exposure questions, and the conventional
counter (``the regulator already certifies it''~\cite{ICNIRP2020})
has not won the argument.  Both deployment blockers route through
\emph{the human body in the cell}.

This note proposes that the same on-device body twin can address
both.  Part~1 turns the conversation by giving each user a
real-time, transparent dose receipt: \emph{this is exactly how much
electromagnetic power your body is currently absorbing.}  Part~2
turns the blockage cost into a small adaptable channel-shaping
opportunity: the body-as-passive-RIS, with pose as the tuning knob
the user already controls.

The technical machinery is in v5
(\texttt{paper\_jsac\_v5.tex}, hereafter \emph{v5}): a per-body
exposure operator $\Qabs^{(u)}$ from a ray-traced path dictionary,
a translation-phasor identity that lets per-slot $\Qabs^{(u)}$
refresh be a diagonal sandwich rather than a fresh surface
integral, an SMPL-X body model with a virtual-IMU pose estimator,
and a per-slot primal projection of regularized ZF that enforces
per-body compliance at the BS side.  The \emph{deployable} part of
v5 sits at the BS.  This note moves the same machinery to the
user's phone and adds two things:
\begin{enumerate}[leftmargin=1.6em,itemsep=2pt]
\item the per-user SAR receipt (\cref{sec:sar-receipt})
\item the differentiability of $\Qabs^{(u)}$ in pose (\cref{sec:diff})
that lets a closed-loop optimisation suggest pose moves to the user.
\end{enumerate}

\subsection{Why ``RIB''?  A literal RIS analogy}

A reconfigurable intelligent surface (RIS) is a passive panel of
sub-wavelength tunable phase shifters that reshapes incident waves
toward a desired receiver.  A human body at mmWave is, in the
high-frequency Kirchhoff/PO regime
(\texttt{rib.tex}~\cref{a:po-rib}--\cref{a:pb-rib}), a passive
scatterer whose surface phase pattern $e^{-ik_0\kout\cdot\rr}$ is
\emph{fixed by pose} and whose only tuning is the user's posture.
The body cannot tune electronically; it can re-pose, on the order of
\SI{1}{\Hz}.  In a one-line slogan from \texttt{NEW\_ANGLE.md}:

\begin{quote}
\emph{Human bodies are reflectors at mmWaves.  They are also mounted
with a highly intelligent organ called the brain.  They are
therefore reflective intelligent human bodies (RIHBs).}
\end{quote}

The cascaded-channel structure familiar from RIS,
$\hh_{\mathrm{cascaded}} = \hh_{\mathrm{BS-to-RIS}}\cdot
\Phimat\cdot\hh_{\mathrm{RIS-to-UE}}$, has a body-side twin
\begin{equation}\label{eq:cascaded-body}
  \hh_{\mathrm{cascaded}}=\hh_{\mathrm{LOS}}+
   \hh_{\mathrm{BS-to-body}}\cdot
   \Frib(\thetab)\cdot
   \hh_{\mathrm{body-to-UE}},
\end{equation}
with $\Frib(\thetab,\kout)$ the per-pose bistatic body scattering
matrix from \texttt{rib.tex}~\cref{def:Fb-rib}.  The body's
``$\Phimat$'' is the implicit phase mask induced by its surface and
pose; the BS's ``re-pointing'' of the body is the pose suggestion
the user takes (or does not take).  Tuning latency is gait-cadence
($\sim$\SI{1}{\Hz}) instead of MHz-cadence; channel-rank lift is
in tens of independent angular modes per body
(\texttt{rib.tex}~\S\,9).  At population scale, even sparse opt-in
fractions translate to a measurable cooperative-RIB lift on
neighbours' channels (cf.\ \texttt{scenarios.md} S4).

We adopt RIS lingo throughout: ``cascaded channel,''
``passive RIS,'' ``surface phase mask,'' ``re-pointing,'' all
re-targeted to the body.  The analogy is informal but pedagogically
useful.

\subsection{Two contributions, one piece of machinery}

The same on-device compute path serves both contributions.
\Cref{tab:two-parts} summarises.

\begin{table}[h]
\caption{The two contributions of this paper, on the same machinery.}
\label{tab:two-parts}
\centering\small
\begin{tabular}{@{}p{0.32\linewidth}p{0.30\linewidth}p{0.30\linewidth}@{}}
\toprule
\textbf{Question} & \textbf{Part~1: SAR receipt} &
\textbf{Part~2: RIHB capacity} \\
\midrule
What is computed? &
$\Pabs(\thetab,\xvec)=\xvec^H\Qabs(\thetab)\xvec$ in real time &
$\nabla_\thetab R(\thetab,\xvec)$ for a pose suggestion $\thetab^*$ \\
Output & a number, displayed in an app & a pose move, suggested via
notification \\
Privacy & receipt computed locally, optionally aggregated &
gradient computed locally, only $\thetab^*$ accept/reject leaves
device \\
Data flow with BS & DL: path dictionary; UL: nothing required &
DL: path dictionary; UL: pose estimate, accept/reject \\
Optimization role & \textbf{none}: informational only &
maximises rate under comfort cost only \\
What v5 already does & nothing: SAR is BS-side compliance only &
nothing: pose is observed, not controlled \\
\bottomrule
\end{tabular}
\end{table}

\subsection{Why the network and the regulator both want this}

The deployment value of Part~1 is consumer-facing.  An app that
shows the user their real-time absorbed power addresses the
trust-deficit conversation directly, in a way that ``it's below
the regulator's limit, take our word for it'' has not.  Several
downstream applications follow.  Per-user logged exposure histories
become input to dose-response epidemiology that the field has
been unable to run for lack of per-user data
(\cite{ICNIRP2020}~\S\,1.4 admits the gap).  Aggregate exposure
maps at city scale become possible from anonymised receipts.  The
regulator gets, as a side effect, the data they have always wanted
but could not collect at scale.

The deployment value of Part~2 is operator-facing.  At binding load
the network already trades \SI{30}{} to \SI{40}{dB} of channel
quality on body-shadowed users; even modest pose moves recover a
fraction of this.  In NLOS scenarios where the user's body is the
only available scatterer, the cascaded-channel
\eqref{eq:cascaded-body} is the only available channel, and the
RIHB-with-pose-control is what makes it predictable.  The
business-model conversation (operators rebating users for opted-in
co-operation, etc.) is left for the discussion.

% ============================================================
\section{Setup and notation}\label{sec:setup}
% ============================================================

The setup is v5's, with the body model promoted from observation to
decision variable.

\subsection{Geometry and array}

A single base station has $M$ antenna elements at known positions
$\{\rr_j\}_{j=1}^M$ and radiates with precoder $\xvec\in\Complex^M$
under power constraint $\|\xvec\|^2\le P_{\mathrm{tx}}$.  In running
examples we follow v5 with $M=64$ ($8\times 8$ uniform rectangular
array, half-wavelength spacing, carrier $f_c=\SI{28}{\GHz}$,
$\lambda\approx\SI{1.07}{\cm}$).

A user is a closed parametric surface
$\Sigma(\thetab,\betab)$ with shape $\betab\in\Reals^{10}$ and pose
$\thetab\in\Reals^{72}$ from SMPL-X (Pavlakos et al., 2019).  The
user holds a phone receiver at $\rr_p(\thetab)$, attached to the
right wrist (an SMPL-X joint), so $\rr_p$ moves smoothly with
$\thetab$.  In the running examples the phone is at
$d_p\sim$\SI{30}{\cm} from the wrist.

\paragraph{Far-field bookkeeping, per anatomical scale.}  The
textbook ``whole-body Fraunhofer distance''
$2D^2/\lambda$ with $D=\SI{1.8}{\m}$ gives \SI{605}{\m} at
\SI{28}{\GHz}.  This number is the conservative bracket for treating
the entire body as one coherent radiating aperture, and is the
formula \texttt{rib.tex}~\cref{a:far-rib} quotes.  It is the wrong
scale for what we actually compute.  No receiver looks at a body as
a single \SI{1.8}{\m} aperture.  The operationally relevant scale
is per-anatomical-patch:
\begin{center}\small
\begin{tabular}{@{}lr@{}}
\toprule
Object & $r_F=2D^2/\lambda$ at \SI{28}{\GHz} \\
\midrule
Whole body, $D=\SI{1.8}{\m}$ & \SI{605}{\m} \\
Torso, $D=\SI{0.5}{\m}$       & \SI{47}{\m} \\
Head, $D=\SI{0.2}{\m}$        & \SI{7.5}{\m} \\
Hand or forearm, $D=\SI{0.1}{\m}$  & \SI{1.9}{\m} \\
Single specular patch, $D=\SI{0.05}{\m}$ & \SI{0.5}{\m} \\
\bottomrule
\end{tabular}
\end{center}
The BS at $\sim\SI{30}{\m}$ and aperture
$\sim\SI{0.14}{\m}$ (8$\times$8 URA, half-wavelength)
has $r_F^{\mathrm{BS}}=\SI{3.7}{\m}$, so the body sits comfortably
in the BS's far field and each BS-to-body path arrives as a clean
plane wave: v5's $\Qabs$ formulation \emph{transfers unchanged}.
The phone at $d_p\sim\SI{30}{\cm}$ from the torso, however, is
inside the torso's $r_F=\SI{47}{\m}$, so any \emph{body-to-phone
bistatic scattering} that involves a torso-scale patch is in the
body's near field and needs the Fresnel-Kirchhoff treatment of
\cref{sec:nearfield}.  Critically, the pose-rate gradient in the
LOS-un-shadowing regime (the hero scenario, PoC at
\texttt{poc/s1\_window/}) does \emph{not} invoke any
body-scattering matrix at all: it is a pure geometric blockage
effect on the direct BS-to-phone LOS, captured by a visibility
function $v(\thetab)\in[0,1]$.  The near-field correction matters
only for the body-as-passive-RIS contribution in NLOS scenarios
(S2, S3); see \cref{rem:nf-applicability}.

\subsection{Path-space factorization (after v5)}

The propagation between the array and any spatial probe is the
$N$-path set $\{j(n),\khat_n,\bm\psi_n\}_{n=1}^N$ from v5~\S\,II,
with $\khat_n$ unit propagation direction at the probe and
$\bm\psi_n$ polarization amplitude.  The body-surface field channel
factors as
\begin{equation}\label{eq:Gtilde}
  \Gtil(\rr;\thetab)=\tilde{\bm\Psi}(\rr;\thetab)\,
                     \Phimat(\rr)\,\Jmat\in\Complex^{3\times M},
\end{equation}
with $\Phimat$ the per-path phase diagonal, $\Jmat$ the
element-to-path dispatch matrix, and $\tilde{\bm\Psi}$ the
tissue-and-Fresnel-weighted per-path read-out (v5~eq.~26).  Per-body
exposure operator and (over the back hemisphere) bistatic operator
are
\begin{align}
  \Qabs(\thetab)
  &= \int_{\Sigma(\thetab)}\Gtil(\rr;\thetab)^H\,\Gtil(\rr;\thetab)\,\diff A(\rr)
  =\Jmat^T\Mmat(\thetab)\Jmat,
  \label{eq:Qabs-def}\\
  \mathbf{Q}_{\mathrm{bi}}(\thetab,\kout)
  &=\frac{1}{Z_0}\,\Frib(\thetab,\kout)^H\,\Frib(\thetab,\kout),
  \label{eq:Qbist-def}
\end{align}
with $\Mmat(\thetab)\in\Complex^{N\times N}$ the per-body path-pair
Gram (v5~eq.~24) and $\Frib$ from \texttt{rib.tex}~\S\,4.

\subsection{The body-to-phone direction, formally}\label{sec:body-phone-dir}

The bistatic operator
$\mathbf{Q}_{\mathrm{bi}}(\thetab,\kout_{u\to\mathrm{UE}})$ above
takes a single outgoing direction $\kout_{u\to\mathrm{UE}}$ as
argument.  In a far-field receiver geometry,
\begin{equation}\label{eq:kout-ff}
  \kout_{u\to\mathrm{UE}}^{\mathrm{ff}}
  \;\equiv\;\frac{\rr_p-\rr_u^{\mathrm{ref}}}
                  {\|\rr_p-\rr_u^{\mathrm{ref}}\|}
  \in S^2,
\end{equation}
with $\rr_u^{\mathrm{ref}}$ a body reference point (centroid is the
canonical choice).  This is the single-direction shorthand that
\texttt{rib.tex} and v5 use throughout.  It is well-defined when
$\|\rr_p-\rr_u^{\mathrm{ref}}\|\gg D_u^2/\lambda$ for $D_u$ the
body's diameter.

\paragraph{Why \eqref{eq:kout-ff} is shorthand at smartphone
distances.}  At $f_c=\SI{28}{\GHz}$, the operationally relevant
scatterer scale is the anatomical patch the path actually
backscatters off (per \cref{sec:setup}'s Fraunhofer table).  For a
torso-scale patch ($D\sim\SI{0.5}{\m}$,
$r_F=\SI{47}{\m}$), the phone at $d_p\sim\SI{30}{\cm}$ is two
orders of magnitude inside the torso's near field.  At this
geometry there is no single $\kout_{u\to\mathrm{UE}}$ for the
torso-scale patch: different points on the body's lit area see
different directions to the phone.  For finger-scale features
($D\sim\SI{0.05}{\m}$, $r_F=\SI{0.5}{\m}$) the phone at \SI{30}{\cm}
is in the patch's far field and the constant-direction shorthand
applies locally.  The right object for the dominant scatterers is
the \emph{per-point outgoing direction}
\begin{empheq}[box=\tcbhighmath]{equation}\label{eq:eta-perpoint}
  \hat{\bm\eta}(\rr;\rr_p)
  \;\equiv\;\frac{\rr_p-\rr}{\|\rr_p-\rr\|}
  \in S^2,\qquad\rr\in\Sigma(\thetab),
\end{empheq}
which varies smoothly across the body surface.  In the Kirchhoff
far-field integral of \texttt{rib.tex} eq.~17, the constant
$\kout$ in the phase factor $e^{-ik_0\kout\cdot\rr}$ is replaced by
$\hat{\bm\eta}(\rr;\rr_p)$ and the $1/(4\pi R)$ prefactor by the
Fresnel-Kirchhoff kernel
$e^{ik_0|\rr-\rr_p|}/(4\pi|\rr-\rr_p|)$.  Both substitutions are
in \texttt{rib.tex}~\S\,12.

\paragraph{Practical convention used in this paper.}  We write
$\kout_{u\to\mathrm{UE}}$ for the conventional far-field direction
\eqref{eq:kout-ff} and use it in informal expressions where the
substitution $\kout\to\hat{\bm\eta}(\rr;\rr_p)$ is understood
to apply pointwise inside the surface integral when the phone is
in the body's near field.  \Cref{sec:nearfield} states the
quantitative consequence (a near-field correction factor
$g_{\mathrm{NF}}(d_p,D)$) for the SAR receipt and the bistatic
gradient.

\paragraph{The BS-to-body direction is far field.}  The BS array
($\sim\SI{0.14}{\m}$ aperture) at $\sim\SI{30}{\m}$ has Fraunhofer
distance $r_F^{\mathrm{BS}}=\SI{3.7}{\m}$, so the body is in the
BS's far field, and per-path directions $\{\khat_n\}$ at the body
are well-defined (no near-field correction needed on the BS side).
The asymmetry matters: the BS sees the body as a small angular
target, but a phone held \SI{30}{\cm} from the body sees the body
as a distributed near-field scatterer at the torso scale.

\begin{remark}[Where the near-field correction actually matters]
\label{rem:nf-applicability}
The pose-rate gradient \eqref{eq:grad-decomp} (introduced below)
has two terms.  The LOS-un-shadowing term is a geometric blockage
modulation of the direct BS-to-phone path; it has no body
scattering matrix and no near-field issue.  The body-mediated
bistatic term involves $\Frib(\thetab,\kout)$ at the
body-to-phone direction; \emph{this} is where the near-field
correction matters when the body part doing the scattering is at a
torso or head scale.  The hero-scenario PoC at
\texttt{poc/s1\_window} measures only the un-shadowing term and is
therefore robust to the near-field issue.  Scenarios S2 (NLOS) and
S3 (multi-user co-operative) need the near-field treatment of
\cref{sec:nearfield}.
\end{remark}

\subsection{The translation-phasor identity (from v5)}\label{sec:translation}

The single most important quantity for both parts is per-slot
$\Qabs(\thetab)$ at the body's current world position
$\rr_0+\mathbf{t}$.  v5~eq.~25 observes that the static
body-and-pose Gram is identity-translatable:
\begin{empheq}[box=\tcbhighmath]{equation}\label{eq:translation}
  \Mmat(\rr_0+\mathbf{t},\thetab)
  =\Phimat(\mathbf{t})^H\Mmat(\rr_0,\thetab)\Phimat(\mathbf{t}),
\end{empheq}
with $\Phimat(\mathbf{t})=\mathrm{diag}(e^{-ik_0\khat_n\cdot
\mathbf{t}})$.  This is what makes both Part~1 (live SAR) and Part~2
(pose gradient) deployable: $\Mmat$ is built once per pose at the
body-local frame ($\sim$\SI{7.5}{\ms} per body on a phone-class GPU
per v5~Table~II), and per-slot translation is an
$O(N)$ phasor sandwich that costs sub-millisecond.  All compute
estimates below assume this trick.

% ============================================================
\section{Part~1: the live SAR receipt}\label{sec:sar-receipt}
% ============================================================

\subsection{What it computes}

The user's phone receives, from the BS, a periodic update of the
local path dictionary
$\{\khat_n,\bm\psi_n,j(n)\}_{n=1}^N$ at the user's location.  At
each slot, the phone computes
\begin{equation}\label{eq:sar-receipt}
  \Pabs(\thetab,\xvec_t)=\xvec_t^H\Qabs(\thetab)\xvec_t
  \;[\mathrm{W}],
\end{equation}
where $\xvec_t$ is the precoder the BS used at the slot (also
broadcast in the path-dictionary update or inferable from the
user's known channel) and $\thetab$ is the user's current
estimated pose from the on-device IMU + AR pipeline.  The receipt
quantity is in absolute Watts at the surface; the user-facing
display can normalize to $\Pabs/L_{\mathrm{RL}}$ (fraction of the
local reference-level budget), or to a body-mass-normalized SAR
($\Pabs/m_{\mathrm{body}}$ in W/kg), or to a 6-min-window time
average for the Brussels regulatory window.

The path dictionary already exists on the BS side for every served
or sensed user (v5~\S\,III.A).  Sharing it down to the user costs
the bandwidth budget of \cref{tab:bandwidth} below
($\sim$\SI{130}{\kbps}), four orders of magnitude below the data
rate.

\subsection{Why this has not been done}

The standard SAR-disclosure mechanism in handsets is
device-specific: a phone's manufacturer publishes a fixed
SAR$_{1g}$ or SAR$_{10g}$ number measured under standardized
worst-case conditions in a specific anatomical phantom.  Real-time
per-user SAR has not been done because (i) it requires real-time
ray tracing of the user's local channel, (ii) it requires a
calibrated body model on-device, and (iii) it requires the BS to
share the path information that real-time SAR depends on.  All
three become available when v5's body twin runs at the BS:
(i) the BS is doing the ray tracing already, (ii) SMPL-X is on
the phone for AR purposes anyway, and (iii) the BS-to-UE channel
already carries DL control plane that can absorb the path
dictionary.

The opportunity-cost argument is therefore:
\emph{the operator builds the body twin for compliance; the user
gets the receipt for free.}

\subsection{Compute architecture}

\begin{center}
\begin{tikzpicture}[
  node distance=10mm and 14mm, font=\small,
  block/.style={draw, rounded corners, align=center, minimum height=8mm,
                minimum width=20mm, fill=gray!8},
  arrow/.style={-{Stealth[length=2mm]}, thick},
]
\node[block] (bs) {BS\\(v5 twin)};
\node[block, right=of bs] (rt) {Path\\dict.\ DL};
\node[block, right=of rt] (uephone) {Phone\\(SMPL-X)};
\node[block, below=of uephone] (qcomp) {$\Qabs(\thetab)$\\via \cref{eq:translation}};
\node[block, below=of qcomp] (display) {SAR\\display};
\node[block, left=of qcomp] (logger) {Logger\\(opt-in)};
\node[block, below=of logger] (epi) {Aggregated\\epi study};
\draw[arrow] (bs)--(rt);
\draw[arrow] (rt)--(uephone);
\draw[arrow] (uephone)--(qcomp);
\draw[arrow] (qcomp)--(display);
\draw[arrow, dashed] (qcomp)--(logger);
\draw[arrow, dashed] (logger)--(epi);
\end{tikzpicture}
\end{center}

The dashed arrows are opt-in.  No body geometry leaves the device;
only the receipt itself flows out, and only with consent.  Aggregated
exposure histories support population-scale studies that the field
has been unable to run for lack of data
(\cite{ICNIRP2020}~Annex~A flags this gap).

\subsection{What the receipt is and what it is not}

\begin{itemize}[leftmargin=1.6em,itemsep=3pt]
\item It \emph{is} a real-time, model-based estimate of the user's
absorbed power from the network's perspective.
\item It \emph{is} a transparent number the user can act on (move,
re-orient the device, defer the session).
\item It \emph{is} a rich data source for population epidemiology
that the standardized SAR$_{1g}$ disclosure cannot supply.
\item It \emph{is not} a regulatory measurement.  Compliance to
ICNIRP basic restrictions is a basic-restriction-pathway
measurement; v5 carries that argument at the BS.
\item It \emph{is not} part of the user-side rate optimisation.
The receipt informs; nothing in the user's gradient trades it.
This is by design: the user is operating in a regime where
compliance is enforced upstream and the receipt is voluntary
self-monitoring.
\end{itemize}

% ============================================================
\section{Part~2: pose-differentiable RIHB capacity}\label{sec:diff}
% ============================================================

The same body twin gives, with one extra step, the gradient
$\nabla_\thetab R(\thetab,\xvec)$ that lets the BS suggest pose
moves to the user.

\subsection{Differentiability of $\Qabs$ in pose}

\paragraph{Smoothness.}  SMPL-X's linear blend skinning makes each
vertex position a $C^\infty$ function of pose almost everywhere
(non-smoothness on a measure-zero set of joint-exchange
configurations).  The triangle Fresnel weights and surface integrand
are $C^\infty$ in vertex positions away from grazing-incidence
silhouettes, where a one-degree-wide GELU smoothing of the
visibility gate (TAP \S\,2.6) restores smoothness with $\le$\SI{0.4}{\percent}
bias on $\tr\Qabs$.  Composing,
$\thetab\mapsto\Qabs(\thetab)$ is $C^\infty$ on the usable pose
manifold, and the chain-rule Jacobian factors as
\begin{equation}\label{eq:Q-jac}
  \frac{\partial\Qabs}{\partial\theta_i}
  =\sum_t\sum_n\sum_{v\in t}
    \frac{\partial q_{t,n}}{\partial\rr_v}
    \cdot\frac{\partial\rr_v}{\partial\theta_i},
\end{equation}
the SMPL-X autograd handing back the $\partial\rr_v/\partial\theta_i$
column.  This is the differentiability statement of v2~\S\,4 in one
remark.  The key consequence is that the pose-conditional rate
$R(\thetab,\xvec)$ has a computable gradient on the device.

\subsection{Cost of one pose-gradient step}

The same v5 cadence table applies.  Per pose-update tick:
\begin{itemize}[leftmargin=1.6em,itemsep=2pt]
\item Build $\Mmat(\thetab)$ in body-local coordinates: \SI{7.5}{\ms}.
\item Per-slot translation phasor on $\Mmat$: sub-\SI{1}{\ms} (v5).
\item One reverse-mode VJP of $R(\thetab,\xvec^*)$ against $\thetab\in
  \Reals^{72}$: $\sim$\SI{15}{\ms} on a phone-class GPU.
\end{itemize}
At the loop cadence of \SI{1}{\Hz} pose suggestions, this is
comfortable.  At slot cadence (\SI{30}{\ms} in the v5 plaza
configuration), the gradient is computed per pose update tick, not
per slot; the precoder runs at slot cadence using the suggested
pose.

\subsection{The rate-only loss}

The user's optimization loss is rate-maximisation under comfort
\begin{empheq}[box=\tcbhighmath]{equation}\label{eq:loss-rate-only}
  \mathcal{L}(\thetab,\xvec;\thetab^{(0)})
  = -R(\thetab,\xvec)
    \;+\;\lambda_M\,C(\thetab-\thetab^{(0)}),
\end{empheq}
where $C$ is a separable comfort cost calibrated such that $C\le 1$
corresponds to a \SI{5}{\degree} torso rotation (RULA score
increment $+1$, ``low load'').  No exposure penalty appears: the
compliance cap is enforced by the BS via v5's projected ZF
(v5~eq.~28) and is not the user's burden to optimize against.  The
SAR receipt of \cref{sec:sar-receipt} is computed and displayed but
does not enter \eqref{eq:loss-rate-only}.

\subsection{The pose gradient}

The rate gradient decomposes naturally into terms that mirror the
RIS cascaded-channel structure of \eqref{eq:cascaded-body}.

\begin{lemma}[Pose gradient decomposition]\label{lem:grad-decomp}
For single-stream MRT precoding $\xvec=\sqrt{P_{\mathrm{tx}}}
\hh(\thetab)/\|\hh(\thetab)\|$ with cascaded channel
$\hh(\thetab)=\hh_{\mathrm{LOS}}(\thetab)+\sum_b\Frib^{(b)}(\thetab)
\xvec/[\cdot]$, below the MCS27 cap,
\begin{equation}\label{eq:grad-decomp}
  \frac{\partial R}{\partial\theta_i}
  =\frac{B}{\ln 2}\,\frac{\SINR_{\max}}{1+\SINR_{\max}\|\hh(\thetab)\|^2}
   \cdot 2\,\RE\bigl(\hh(\thetab)^H\,\partial_{\theta_i}\hh(\thetab)\bigr),
\end{equation}
with $\partial\hh/\partial\theta_i=\partial\hh_{\mathrm{LOS}}/
\partial\theta_i+\sum_b\partial\hh_{\mathrm{body},b}/\partial\theta_i$.
The first term is the LOS un-shadowing gradient; the second is the
RIB-mediated bistatic re-routing gradient.
\end{lemma}

\noindent The two terms have different magnitudes in different
geometries.  In the open-window LOS case (\texttt{scenarios.md}~S1),
the un-shadowing term dominates and dictates the gradient
direction.  In the NLOS case (S2), the LOS term is identically zero
and the RIB-mediated term is what the gradient sees.  In the
multi-user co-operative case (S3), each user's body shows up in
the others' RIB-mediated term; the gradient is coupled and the
joint optimum is a correlated equilibrium.

\subsection{What pose moves do at realistic geometries}

The PoC \texttt{poc/s1\_window/findings.md} measured the
single-DOF (torso yaw) sweep in a binding-regime body-shadowed LOS
scenario at $f_c=\SI{28}{\GHz}$, $M=64$ URA, BS at \SI{30}{\m},
\SI{35}{\dB} scene loss.  Headline: a \SI{15}{\degree} torso turn
buys $+\SI{14.5}{\percent}$ rate ($+\SI{2.6}{\dB}$ SINR) at no
operationally-relevant exposure cost.  The pre-experiment estimate
of $+\SI{10}{\dB}$ to $+\SI{18}{\dB}$ rate was wrong: at MCS27
saturation, dB rate gain compresses against the cap, and the
honest hero number is in absolute Mbps or in pre-cap SINR dB.

\begin{remark}[Body-as-RIS gain ceiling]\label{rem:ceiling}
The body's projected area is $\sim\SI{0.10}{\m^2}$ to
$\SI{0.15}{\m^2}$ in any single direction.  At the array's
\SI{30}{\m}-distance Fresnel zone of $\sim\SI{40}{\cm}$ radius, the
body covers $\sim 30\,\%$ of the Fresnel zone area.  Pose moves
that swap which $30\,\%$ of the Fresnel zone is body-covered shift
the SINR by at most the $\sim\SI{5}{\dB}$ implied by going from
$30\,\%$-blocked to fully-clear.  Below the MCS cap this is a
real $\SI{5}{\dB}$ rate gain; above the cap it manifests as
slack-on-cap headroom and becomes an absolute-Mbps gain only via
re-paired served-user counts.  In NLOS settings the ceiling
depends instead on the cascaded RIB gain $\|\Frib\|$, which scales
as $(1-T_0)A_{\mathrm{lit}}$; at \SI{28}{\GHz} this is
$\sim\SI{0.06}{\m^2}$ per body.
\end{remark}

% ============================================================
\section{Closed loop}\label{sec:loop}
% ============================================================

\subsection{Block diagram}

\begin{center}
\begin{tikzpicture}[
  node distance=10mm and 14mm, font=\small,
  block/.style={draw, rounded corners, align=center, minimum height=9mm,
                minimum width=22mm, fill=gray!8},
  arrow/.style={-{Stealth[length=2mm]}, thick},
]
\node[block] (bs) {BS array};
\node[block, right=of bs] (rt) {Ray tracer\\(v5 twin)};
\node[block, right=of rt] (channel) {Path dict.\\DL channel};
\node[block, below=of channel] (uephone) {Phone twin\\(SMPL-X)};
\node[block, below=of bs] (precoder) {Precoder\\(proj.\ ZF)};
\node[block, below=of precoder] (comp) {Compliance\\enforcer (v5)};
\node[block, right=of uephone, xshift=10mm] (loss) {Rate loss\\$+$ comfort};
\node[block, right=of loss] (suggest) {Pose\\suggest};
\node[block, below=of uephone] (sar) {SAR receipt\\display};
\draw[arrow] (rt)--(channel);
\draw[arrow] (channel)--(uephone);
\draw[arrow] (uephone)--(loss);
\draw[arrow] (loss)--(suggest);
\draw[arrow] (suggest.south) -- ++(0,-7mm) -| (rt.south);
\draw[arrow] (uephone)--(sar);
\draw[arrow] (rt)--(precoder);
\draw[arrow] (precoder)--(comp);
\draw[arrow] (comp)--(bs);
\draw[arrow, dashed] (suggest) -- ++(0,15mm) -| (uephone.north);
\end{tikzpicture}
\end{center}

\noindent v5 owns the BS-side compliance loop (`precoder $\to$
compliance enforcer $\to$ BS array').  This paper adds the
UE-side capacity loop (`channel $\to$ phone twin $\to$ rate loss
$\to$ pose suggest') and the parallel SAR receipt branch.  The two
sides share the path dictionary as the only common data flow.

\subsection{Bandwidth budget}\label{sec:bandwidth}

Each path is described by 8 floats (2 for $\khat\in S^2$, 6 for
$\bm\psi\in\Complex^3$ normalised, 1 each for delay and amplitude).
At $N=50$ paths re-evaluated at \SI{10}{\Hz} and 32-bit precision:

\begin{table}[h]
\centering\small
\caption{Bandwidth budget for the closed loop.}
\label{tab:bandwidth}
\begin{tabular}{@{}lll@{}}
\toprule
Channel & Payload & Rate \\
\midrule
DL: path dictionary           & $50\times 8$ floats $\times$ \SI{10}{\Hz} & \SI{128}{\kbps} \\
DL: pose suggestion (BS-side) & 72 floats $\times$ \SI{1}{\Hz}            & \SI{2.3}{\kbps} \\
UL: pose estimate $\hat\thetab$ & 72 floats $\times$ \SI{1}{\Hz}          & \SI{2.3}{\kbps} \\
UL: accept/reject + SAR opt-in & one-bit + occasional bits                & negligible \\
UL: shape $\betab$ (one-time)  & 10 floats once per session              & \SI{40}{\B} \\
\bottomrule
\end{tabular}
\end{table}

\noindent Total bidirectional control bandwidth $\le$\SI{135}{\kbps},
four orders of magnitude below the data rate the loop negotiates
over.  This is the same scale as the v5 BS-side cadence
(v5~Table~II).

\subsection{Privacy delta against v5}

In v5, the user's pose estimate flows to the BS through the
operator's application channel for compliance computation
(v5~\S\,III.B); body geometry $\betab$ is implicit in the SMPL-X
model the BS already uses to compute $\Qabs^{(u)}$.  This paper's
addition is that the \emph{pose-suggestion derivation} runs
on-device, so only the suggested move and the accept/reject bit
flow back to the BS.  The user's own SAR receipt is computed
on-device and never leaves unless the user opts in to aggregation.
The privacy delta over v5 is therefore: same input data flow,
plus on-device gradient compute, plus opt-in receipt aggregation.

% ============================================================
\section{Connection to v5 and to the JSAC2 SI}\label{sec:v5-and-si}
% ============================================================

\subsection{v5 is the BS-side compliance paper; this is the UE-side
counterpart}

v5's body twin runs on the BS, observes pose at the v5 \SI{1}{\s}
cadence, and uses $\Qabs^{(u)}$ to set the projection scale
$\alpha=\sqrt{\eta_S\,\min_u L^{(u)}/\PabsZF^{(u)}}$ that holds
every body's absorbed power below the cap.  The body twin is
load-bearing on the BS for compliance.  This paper adds two new
on-device uses of the same machinery: live SAR receipts and a
pose-gradient outer loop.  Neither competes with v5; both extend
it.

The empirical findings of v5 carry over.  In particular:
\begin{itemize}[leftmargin=1.6em,itemsep=2pt]
\item v5~\S\,VI.A: $\Qabs^{(u)}$ has effective rank 3 to 4 at the
$99\,\%$ trace fraction, set by the panel's $\sim\SI{14}{\degree}$
angular resolution.  The pose-gradient $\nabla_\thetab Q$ inherits
this rank, so the achievable channel-shaping by pose is in $\sim 4$
dominant modes per body.  Translation: pose moves are
low-dimensional levers on a low-rank operator.
\item v5~Appendix~C: the unconditional projected-area Cauchy bound
on $\xvec^H\Qabs^{(u)}\xvec$ is breached by $49\,\%$ to $70\,\%$ of
unconstrained precoders with median ratio up to $+\SI{6}{\dB}$.
The conditional bound is the right one.  For the SAR receipt this
matters because a no-twin upper bound would systematically
under-estimate the user's actual absorbed power; the receipt has
to use the actual $\Qabs$, not a worst-case envelope.
\item v5~\S\,VI.D: pose accuracy modulates compliance at
sub-percent scale.  For \emph{capacity} (this paper, not v5), pose
accuracy modulates the rate gradient at the level of the
gradient's geometric scale, $\sim\SI{0.7}{\dB}$ per joint-degree
RMS at the PoC's geometry; pose precision matters more for
capacity than for compliance.
\end{itemize}

\subsection{What this paper does not relitigate}

\begin{itemize}[leftmargin=1.6em,itemsep=2pt]
\item Path dictionary construction and CSI-amplitude calibration
(v5~\S\,III.E).
\item Per-slot primal projection of regularised ZF
(v5~\cref{eq:primal-projection}).
\item The dual-ascent QCQP solver
(v5~\cref{app:dualascent}).
\item Tier-A/B/C/D telemetry classification (v5~\S\,III.B).
\item The ISAC bystander detection link budget (v5~\S\,VI.E).
\end{itemize}

\subsection{Mapping to JSAC2 SI scope}

The SI keywords map cleanly: \emph{closed-loop optimization}
(\cref{sec:loop}), \emph{differentiable control} (\cref{sec:diff}),
\emph{twin-in-the-loop} (Part~1 \emph{is} the twin-in-the-loop on
the user side), \emph{cross-layer} (PHY rate $\times$ biomechanical
pose), \emph{application-aware} (the application is the user's
session and the user's exposure-receipt preference), \emph{trust
and transparency} (Part~1 \emph{is} the trust artifact).

% ============================================================
\section{Near-field correction for body-to-phone}\label{sec:nearfield}
% ============================================================

The phone is at $d_p\sim$\SI{30}{\cm} from the body centroid; the
body's Fraunhofer distance at \SI{28}{\GHz} is
$D^2/\lambda\sim\SI{300}{\m}$, so the phone is in the body's deep
near field.  The body scattering matrix $\Frib$ from
\texttt{rib.tex}~\S\,4 is a far-field object; its near-field
analog from \texttt{rib.tex}~\S\,12 replaces the $1/(4\pi R)$
kernel with the Fresnel-Kirchhoff
$e^{ik_0|\rr-\rr_p|}/(4\pi|\rr-\rr_p|)$.  The Pareto bound from
\texttt{rib.tex}~\S\,8 develops a near-field correction factor
$g_{\mathrm{NF}}(d_p,D)\le 2$ for the worst-case torso-to-hand
geometry, which modifies the body-universal $(1-T_0)/T_0$ slope by
at most a factor of two.  For the SAR receipt this matters
quantitatively (the displayed dose is right within a factor of
$g_{\mathrm{NF}}$ at user-and-phone geometries, which we assume
calibrated against full-wave reference simulation per device); for
the capacity gradient this is in the bistatic-scattering term and
is small in the LOS-un-shadowing-dominated regime of the hero
scenario.

% ============================================================
\section{What this is not}\label{sec:notthis}
% ============================================================

\begin{enumerate}[leftmargin=1.6em,itemsep=3pt]
\item Not a compliance paper.  v5 owns BS-side compliance.
\item Not a beamforming paper.  The precoder is v5's projected ZF.
\item Not a pose-estimation paper.  The pose estimator is the
on-device IMU + AR pipeline (Madgwick / Mahony / SMPL-X-from-IMU).
\item Not a JCAS paper.  We use on-device pose tracking, not
backscatter pose estimation.
\item Not a regulatory measurement.  The SAR receipt is a
model-based estimate; the basic-restriction-pathway audit is what
the regulator certifies.
\item Not new hardware.  The BS is a vanilla 5G FR2 panel; the UE
is a vanilla smartphone with IMU and AR.
\item Not a literal RIS paper.  The body cannot tune electronically;
the analogy is geometric (passive surface, gait-tunable phase mask)
and is used pedagogically.
\end{enumerate}

% ============================================================
\section{Loose threads and conjectures}\label{sec:threads}
% ============================================================

\begin{enumerate}[leftmargin=1.6em,itemsep=3pt]
\item \emph{Multi-user RIHB.}  Joint loss
$\mathcal{L}(\thetab_1,\thetab_2,\xvec)$.  Conjecture: a
Stackelberg-style lift exists when both users follow gradients.
Scenario S3 measures.
\item \emph{The discussion observation.}  Re-routing multipath
through one's own body via RIHB pose moves probably \emph{increases}
that user's own absorbed dose (more incident power couples into
absorbing tissue).  This is an interesting observation for the
discussion; it does not feed back into the optimization because
exposure is informational only.  The receipt displays it; the user
decides.
\item \emph{Pose-conditional dose-effect studies.}  Aggregated
SAR receipts at population scale support exposure-dose-response
epidemiology that the field has been unable to run.  The
study-design question is its own paper.
\item \emph{Adversarial poses for stress-testing the BS.}  Worst-case
pose on the comfort ball gives the BS a regret bound on rate-loss
under non-cooperative pose.  An interesting characterization of
v5's robustness to its own pose-estimator quality.
\item \emph{Federated $\betab$.}  Operator estimates population
$\bar\betab$ via federated averaging without seeing any individual
$\betab$; whether the smoothed-mannequin shape is operationally
useful is unknown.
\item \emph{Sub-\SI{6}{\GHz}.}  Pseudo-Brewster collapse holds at
mmWave; below \SI{6}{\GHz} the body's $T_0$ is
polarization-dependent and the body-RIS analogy weakens.
\item \emph{Comfort metric calibration.}  Cost weights $c_i$
calibrated against a small ergonomic user study (RULA / REBA based)
becomes a separate companion study.
\item \emph{Privacy preservation of the receipt aggregation.}
Differential-privacy treatments of the per-slot SAR sample;
secure-aggregation protocols.  Out of scope for this paper.
\item \emph{The nullable Cauchy bound.}  v5~Appendix~C shows the
unconditional bound is loose; the conditional bound holds.  An open
question: what is the equivalent statement for the bistatic
operator, and does it bound the gradient?
\item \emph{The two-paper option.}  Robin's note suggests Part~1
and Part~2 may merit separate papers (the SAR app to a
dosimetry-leaning venue, the RIHB capacity to a comms-leaning one).
The unified spine (this paper) keeps both contributions on one
piece of machinery; splitting is mechanical.
\end{enumerate}

% ============================================================
\section{Summary}\label{sec:sum}
% ============================================================

A user with a smartphone is at mmWave a slowly tunable RIS, and a
real-time exposure data source.  v5 already gives every
infrastructure piece needed to run both halves: a per-body
exposure operator, a translation phasor identity for fast per-slot
refresh, an SMPL-X body model on-device for AR purposes, and a
calibrated path dictionary at the BS.  This paper takes the same
machinery, runs it on the phone, and offers two complementary
products.  Live SAR receipts give the user transparent access to
their own absorbed power, addressing the population-trust
deficit that has slowed mmWave adoption.  The differentiable
pose-gradient on the same operator suggests pose moves that
recover capacity in body-shadowed line-of-sight scenarios and
opens NLOS scenarios via the body-as-passive-RIS effect.  Neither
contribution competes with v5's BS-side compliance; both extend it
with no new hardware on either side, and at a control-loop
bandwidth four orders of magnitude below the data rate the loop
negotiates over.

\end{document}
